\chapter{Data Preprocessing and Dimensionality Reduction}
\label{chap:preprocessing_and_dimensionality_reduction}

The algorithms discussed in the preceding chapters, from the geometric purity of Support Vector Machines to the probabilistic rigor of Bayesian classifiers, all share a common, unspoken assumption: they expect to be fed clean, complete, and well-behaved numerical data. The raw data harvested from the real world, however, is rarely so obliging. It is often incomplete, inconsistent, and riddled with a mixture of numerical and categorical formats. Furthermore, modern datasets are frequently characterized by an immense number of features, a condition that can paradoxically degrade model performance through a phenomenon known as the "curse of dimensionality."

This chapter addresses these pragmatic challenges head-on. It is dedicated to the crucial, often unglamorous, but fundamentally important stages of the machine learning workflow that occur before a single model is trained. We will first establish a systematic workflow for data preprocessing, tackling the essential tasks of handling missing values, encoding categorical variables, and scaling numerical features. Subsequently, we will embark on a deep dive into the theory and application of dimensionality reduction, exploring both feature selection and feature extraction. We will derive from first principles the two most important techniques in this domain: the unsupervised method of Principal Component Analysis (PCA) for maximizing variance, and the supervised method of Linear Discriminant Analysis (LDA) for maximizing class separability. This chapter is the bridge from pristine theory to the messy, complex, and ultimately more rewarding reality of applied data science.

\section{The Data Preprocessing Workflow}
\label{sec:preprocessing_workflow}
Data preprocessing is a series of transformations applied to raw data to make it suitable for machine learning algorithms. A robust and thoughtful preprocessing pipeline is often more impactful on a model's final performance than the choice of the algorithm itself.

\subsection{Handling Missing Data}
\label{subsec:missing_data}
Missing values, often represented as \texttt{NaN} (Not a Number), are a ubiquitous problem. Their presence can cause algorithms to fail or, more insidiously, can introduce significant bias into the model if not handled correctly.

\subsubsection*{Strategies for Elimination}
The simplest approach is to remove samples or features that contain missing values.
\begin{itemize}
    \item \textbf{Listwise Deletion}: Any row (sample) containing at least one missing value is discarded. This is a viable option only if the number of affected rows is very small (e.g., $<5\%$ of the dataset) and the missingness is completely random. If the missingness is systematic (e.g., a certain demographic group is less likely to answer a survey question), this method will introduce bias into the training data.
    \item \textbf{Column Deletion}: If a particular feature (column) has a very high percentage of missing values (e.g., $>50\%$), it may carry little useful information. In such cases, it might be better to drop the entire feature column.
\end{itemize}

\subsubsection*{Strategies for Imputation}
Imputation is the process of filling in, or "imputing," missing values with a substitute.
\begin{itemize}
    \item \textbf{Mean/Median Imputation}: This is the most common technique. The missing values in a numerical feature are replaced with the mean or median of the non-missing values in that same feature. The \textbf{median} is generally preferred as it is robust to outliers, whereas the mean can be skewed by extreme values.
    \item \textbf{Mode Imputation}: For categorical features, the most logical substitute is the most frequently occurring category, i.e., the \textbf{mode}.
    \item \textbf{More Advanced Techniques}: For greater accuracy, missing values can be treated as a prediction problem in themselves. Methods like \textbf{K-Nearest Neighbors (KNN) Imputation} use the values of the $k$ most similar complete samples to impute the missing value. More complex methods can even train a regression model to predict the missing value based on other features.
\end{itemize}

\subsection{Manipulating Categorical Data}
\label{subsec:categorical_data}
Machine learning models are mathematical functions that require numerical input. Therefore, categorical data (text labels) must be converted into a numerical representation.

\subsubsection*{Ordinal vs. Nominal Features}
It is crucial to first distinguish between two types of categorical data:
\begin{itemize}
    \item \textbf{Ordinal Features}: These are categories that have an intrinsic, natural order (e.g., `['low', 'medium', 'high']` or `['cold', 'warm', 'hot']`).
    \item \textbf{Nominal Features}: These are categories that do not have any inherent order (e.g., `['red', 'green', 'blue']` or `['USA', 'France', 'Japan']`).
\end{itemize}

\subsubsection*{Encoding Techniques}
\begin{itemize}
    \item \textbf{Label Encoding (for Ordinal Features)}: For ordinal data, a simple mapping of each category to an integer is appropriate, as it preserves the inherent order. For example: `{'low': 1, 'medium': 2, 'high': 3}`.
    \item \textbf{One-Hot Encoding (for Nominal Features)}: Applying simple label encoding to nominal data is a critical error, as it introduces an artificial and misleading order that the model may misinterpret (e.g., implying that `Japan > France > USA`). The correct technique is \textbf{One-Hot Encoding}. This method creates a new binary (0 or 1) feature for each unique category. For a given sample, the column corresponding to its category will have a value of 1, and all other new columns will be 0.
\end{itemize}
\begin{remark}[The Dummy Variable Trap]
When using one-hot encoding, if a feature has $k$ unique categories, the resulting $k$ new columns will be perfectly multicollinear (i.e., one column can be perfectly predicted from the others). For many linear models, this can cause issues with model fitting. The standard practice is to drop one of the $k$ new columns, leaving $k-1$ columns that contain all the necessary information without the multicollinearity. This is known as the dummy variable trap.
\end{remark}

\subsection{Feature Scaling}
\label{subsec:feature_scaling}
Feature scaling is a critical step for algorithms that are sensitive to the magnitude of input features. This includes all algorithms that compute distances (like KNN and SVMs) or use gradient descent for optimization (like Linear/Logistic Regression and neural networks). If one feature ranges from 0 to 1,000,000 and another ranges from 0 to 1, the former will completely dominate the model's learning process.

\subsubsection{Normalization as an Affine Transformation (Min-Max Scaling)}
Normalization rescales features to a fixed range, typically $[0, 1]$. This is not merely a formula, but a specific case of an \textbf{affine transformation}, which has the general form $x' = ax + b$. Our goal is to find the coefficients $a$ and $b$ that map the original domain $[x_{\min}, x_{\max}]$ to the target range $[0, 1]$.

We can derive these coefficients by imposing two boundary conditions:
\begin{enumerate}
    \item When an input value is the minimum of its feature, $x = x_{\min}$, its transformed value must be the minimum of the new range, $x'_{\text{norm}} = 0$.
    \item When an input value is the maximum, $x = x_{\max}$, its transformed value must be the maximum, $x'_{\text{norm}} = 1$.
\end{enumerate}
This gives us a system of two linear equations with two unknowns, $a$ and $b$:
\begin{align*}
    a \cdot x_{\min} + b &= 0 \quad &(1) \\
    a \cdot x_{\max} + b &= 1 \quad &(2)
\end{align*}
Subtracting equation (1) from equation (2) allows us to solve for $a$:
$$ a \cdot x_{\max} - a \cdot x_{\min} = 1 \implies a(x_{\max} - x_{\min}) = 1 \implies a = \frac{1}{x_{\max} - x_{\min}} $$
Substituting this result for $a$ back into equation (1) allows us to solve for $b$:
$$ \frac{x_{\min}}{x_{\max} - x_{\min}} + b = 0 \implies b = - \frac{x_{\min}}{x_{\max} - x_{\min}} $$
Now, substituting our derived coefficients $a$ and $b$ back into the general affine transformation form $x'_{\text{norm}} = ax + b$, we naturally arrive at the familiar formula:
$$ x'_{\text{norm}} = \left(\frac{1}{x_{\max} - x_{\min}}\right)x - \frac{x_{\min}}{x_{\max} - x_{\min}} = \frac{x - x_{\min}}{x_{\max} - x_{\min}} $$
This derivation clarifies that Min-Max scaling is a linear stretching and shifting of the data's original distribution to fit a new, bounded interval. Its primary drawback remains its high sensitivity to outliers, as the values of $x_{\min}$ and $x_{\max}$ are dictated by the most extreme points in the dataset.

\subsubsection{Standardization as Statistical Normalization (Z-score Scaling)}
Standardization is the most common and generally the most robust scaling technique. Its goal is not to bind data to a specific range, but to transform it such that it has the properties of a \textbf{standard normal distribution}: a mean ($\mu$) of 0 and a standard deviation ($\sigma$) of 1. This transformation is fundamental in statistics and is known as computing the \textbf{Z-score}.
$$ x_{\text{std}}^{(i)} = \frac{x^{(i)} - \mu}{\sigma} $$
The value of the Z-score has a powerful statistical interpretation: it represents the number of standard deviations a given data point $x^{(i)}$ is away from the mean of its distribution. A Z-score of 2.5 means the point is 2.5 standard deviations above the mean.

While this transformation can be applied to any feature regardless of its underlying distribution, it is particularly meaningful when the feature is approximately normally distributed. In this case, the Z-score directly relates to the probability of observing such a value. This connects to the empirical 68-95-99.7 rule: approximately 95\% of the data will have a Z-score between -2 and 2, and values with a Z-score greater than 3 are exceptionally rare. This property makes the Z-score a cornerstone of statistical hypothesis testing and outlier detection. Because standardization uses the mean and standard deviation, it is less sensitive to extreme outliers than Min-Max scaling and is often the default choice for preparing data for a wide range of machine learning algorithms.

\section{Dimensionality Reduction}
\label{sec:dimensionality_reduction}
While adding more features can sometimes provide more information to a model, there is a point of diminishing returns. In very high-dimensional spaces, data becomes extremely sparse, and the performance of many algorithms can degrade. This phenomenon is known as the \textbf{curse of dimensionality}. Dimensionality reduction techniques aim to mitigate this problem by reducing the number of features while retaining as much of the relevant information as possible.

\subsection*{Feature Selection vs. Feature Extraction}
There are two main approaches to dimensionality reduction:
\begin{itemize}
    \item \textbf{Feature Selection}: We select a subset of the original features and discard the rest. The goal is to find the most informative subset.
    \item \textbf{Feature Extraction}: We create a new, smaller set of features by combining the original ones. These new features are projections or combinations of the original feature set.
\end{itemize}

\subsection*{Feature Selection: Sequential Backward Selection (SBS)}
Sequential Backward Selection is a classic example of a "wrapper" method for feature selection. It is a greedy, iterative algorithm that aims to find the optimal feature subset of size $k$.
The algorithm proceeds as follows:
\begin{enumerate}
    \item Start with the full set of features, $D$, with size $N$.
    \item Define a criterion function $J$ to evaluate model performance (e.g., accuracy).
    \item At each step, temporarily remove each feature $x^-$ from the current feature set and evaluate the performance $J(D - x^-)$.
    \item Permanently remove the feature that results in the smallest decrease (or largest increase) in performance.
    \item Repeat step 3 and 4 until the feature set is of the desired size, $k$.
\end{enumerate}

\subsection*{Feature Extraction: Principal Component Analysis (PCA)}
Principal Component Analysis is the most widely used unsupervised algorithm for feature extraction. Its goal is to find a lower-dimensional representation of the data that captures the maximum possible variance. It achieves this by identifying a new set of orthogonal axes, called \textbf{principal components}, that align with the directions of highest variance in the data.

\subsubsection*{The Mathematical Derivation of PCA}
The derivation of PCA is an exercise in linear algebra, rooted in the eigendecomposition of the covariance matrix.
\begin{enumerate}
    \item \textbf{Standardize the Data}: Let $X$ be the $m \times N$ data matrix, where $m$ is the number of samples and $N$ is the number of features. It is crucial to first standardize each feature to have zero mean and unit variance.
    \item \textbf{Compute the Covariance Matrix}: Calculate the $N \times N$ covariance matrix, $\Sigma$, of the standardized features.
    $$ \Sigma = \frac{1}{m-1} X^T X $$
    \item \textbf{Perform Eigendecomposition}: The covariance matrix is symmetric, which guarantees that its eigenvectors are orthogonal. We solve the eigenvalue equation:
    $$ \Sigma v = \lambda v $$
    This yields $N$ eigenvectors $v$ and their corresponding eigenvalues $\lambda$. The \textbf{eigenvectors} represent the directions of the principal components, and the \textbf{eigenvalues} represent the magnitude of the variance along those directions.
    \item \textbf{Sort Components}: Sort the eigenvalue-eigenvector pairs in descending order based on the eigenvalues. The eigenvector with the largest eigenvalue is the first principal component, capturing the most variance in the data.
    \item \textbf{Construct the Projection Matrix}: To reduce the dimensionality from $N$ to $k$ (where $k < N$), we select the top $k$ eigenvectors and form an $N \times k$ projection matrix, $W$.
    \item \textbf{Transform the Data}: The new, lower-dimensional feature matrix, $X'$, is obtained by projecting the original standardized data onto the new subspace defined by $W$:
    $$ X' = XW $$
\end{enumerate}
The resulting matrix $X'$ has dimensions $m \times k$. The new features are linear combinations of the old ones and are mutually uncorrelated.

\subsection*{Feature Extraction: Linear Discriminant Analysis (LDA)}
While PCA is a powerful tool, it is an unsupervised method that is blind to the class labels of the data. It simply finds the directions of maximum variance. \textbf{Linear Discriminant Analysis (LDA)}, in contrast, is a \textbf{supervised} dimensionality reduction technique. Its objective is not to maximize variance, but to find the feature subspace that \textbf{maximizes the separability between the different classes}.

\subsubsection*{The Mathematical Intuition of LDA}
LDA seeks to find a projection that maximizes the ratio of the "between-class scatter" to the "within-class scatter." In other words, it wants to project the data such that the means of the different classes are as far apart as possible, while the data points within each class are as close together as possible.

\subsubsection*{The Formal Derivation}
The algorithm involves the following steps:
\begin{enumerate}
    \item Standardize the data.
    \item For each class $k$, compute the mean vector $m_k$.
    \item Compute the \textbf{within-class scatter matrix}, $S_W$, which is the sum of the covariance matrices for each class:
    $$ S_W = \sum_{k=1}^{K} S_k = \sum_{k=1}^{K} \sum_{x \in \alpha_k} (x - m_k)(x - m_k)^T $$
    \item Compute the \textbf{between-class scatter matrix}, $S_B$, which measures the scatter of the class means around the overall mean $m$:
    $$ S_B = \sum_{k=1}^{K} N_k (m_k - m)(m_k - m)^T $$
    \item Solve the generalized eigenvalue problem for the matrix $S_W^{-1}S_B$:
    $$ S_W^{-1}S_B v = \lambda v $$
    \item The eigenvectors $v$ corresponding to the largest eigenvalues are the linear discriminants. We form a projection matrix $W$ with the top $k$ eigenvectors (where $k \le K-1$).
    \item The data is transformed using $X' = XW$, just as in PCA.
\end{enumerate}

\subsection*{Non-Linear Dimensionality Reduction: Kernel PCA (KPCA)}
Both PCA and LDA are linear techniques; they can only find linear subspaces. For datasets with complex, non-linear structures (like spirals or concentric circles), they will fail. Kernel PCA extends PCA to handle such problems by using the kernel trick we introduced in the SVM chapter.

The conceptual steps are:
\begin{enumerate}
    \item Choose a kernel function, $K(x_i, x_j)$.
    \item Use the kernel function to compute the Gram matrix, $K$, for the dataset, where $K_{ij} = K(x_i, x_j)$. This matrix implicitly represents the dot products of the data points in a high-dimensional feature space.
    \item Center the kernel matrix.
    \item Perform eigendecomposition on the centered kernel matrix. The eigenvectors of this matrix will give the projections of the data onto the principal components in the high-dimensional space.
\end{enumerate}
KPCA allows for the discovery of non-linear manifolds in the data, providing a much more powerful tool for feature extraction when linearity cannot be assumed.

\subsubsection{Kernel Principal Component Analysis (KPCA)}
\paragraph{Reformulating PCA via the Gram Matrix.}
The standard PCA algorithm requires the eigendecomposition of the $N \times N$ covariance matrix $\mathbf{\Sigma} = \frac{1}{m}\mathbf{X}^T\mathbf{X}$. Let the eigenvectors be $\mathbf{v}$ and eigenvalues be $\lambda$, such that $\mathbf{\Sigma}\mathbf{v} = \lambda\mathbf{v}$. We can show that any eigenvector $\mathbf{v}$ can be expressed as a linear combination of the input data samples (which form a basis for the feature space):
$$ \mathbf{v} = \sum_{i=1}^m \alpha_i \mathbf{x}_i = \mathbf{X}^T \boldsymbol{\alpha} $$
where $\boldsymbol{\alpha}$ is an $m \times 1$ vector of coefficients. Substituting this into the eigenvalue equation:
$$ \left(\frac{1}{m}\mathbf{X}^T\mathbf{X}\right)(\mathbf{X}^T\boldsymbol{\alpha}) = \lambda (\mathbf{X}^T\boldsymbol{\alpha}) $$
Left-multiplying by $\mathbf{X}$:
$$ \frac{1}{m}\mathbf{X}\mathbf{X}^T\mathbf{X}\mathbf{X}^T\boldsymbol{\alpha} = \lambda \mathbf{X}\mathbf{X}^T\boldsymbol{\alpha} $$
Let us define the $m \times m$ \textbf{Gram matrix} as $\mathbf{K} = \mathbf{X}\mathbf{X}^T$. Each element $K_{ij} = \mathbf{x}_i^T \mathbf{x}_j$ is the dot product between two samples. Substituting $\mathbf{K}$, the equation becomes:
$$ \frac{1}{m}\mathbf{K}\mathbf{K}\boldsymbol{\alpha} = \lambda \mathbf{K}\boldsymbol{\alpha} \implies \mathbf{K}\boldsymbol{\alpha} = m\lambda\boldsymbol{\alpha} $$
This is a standard eigenvalue problem for the Gram matrix $\mathbf{K}$. By solving for the eigenvectors $\boldsymbol{\alpha}$ of $\mathbf{K}$, we can recover the original principal components $\mathbf{v}$. This reformulation shows that PCA can be performed entirely using the dot products between samples, without ever explicitly forming the covariance matrix. This is the key that unlocks the kernel trick.

\paragraph{The KPCA Algorithm and Kernel Centering.}
The KPCA algorithm applies this reformulated PCA to the data in a high-dimensional feature space $\mathcal{F}$ via a mapping $\phi$. The algorithm becomes:
\begin{enumerate}
    \item Construct the kernel matrix $\mathbf{K}$, where $K_{ij} = k(\mathbf{x}_i, \mathbf{x}_j) = \phi(\mathbf{x}_i)^T \phi(\mathbf{x}_j)$.
    \item Center the kernel matrix in the feature space.
    \item Perform eigendecomposition on the centered kernel matrix to find its eigenvectors $\boldsymbol{\alpha}$, which define the principal components.
\end{enumerate}
The centering step is non-trivial. PCA requires the data to be centered. In the feature space $\mathcal{F}$, this means we need to work with $\phi_c(\mathbf{x}_i) = \phi(\mathbf{x}_i) - \frac{1}{m}\sum_{l=1}^m \phi(\mathbf{x}_l)$. The centered kernel matrix $\mathbf{K}'$ would have elements $K'_{ij} = \phi_c(\mathbf{x}_i)^T \phi_c(\mathbf{x}_j)$. Instead of computing this explicitly, we can derive a formula to compute $\mathbf{K}'$ directly from $\mathbf{K}$. By expanding the dot product and expressing the mean terms as averages over the rows and columns of the original kernel matrix $\mathbf{K}$, we arrive at the centering formula:
$$ \mathbf{K}' = \mathbf{K} - \mathbf{1}_m\mathbf{K} - \mathbf{K}\mathbf{1}_m + \mathbf{1}_m\mathbf{K}\mathbf{1}_m $$
where $\mathbf{1}_m$ is an $m \times m$ matrix where every element has the value $1/m$. This allows us to perform the entire PCA procedure in the high-dimensional feature space without ever knowing the coordinates of the mapped points, thus completing the kernel trick.

\part{Scalable Computing and Big Data Paradigms}

The data processing and feature engineering techniques discussed thus far implicitly assume that the entire dataset can be held and manipulated within the memory of a single machine. However, the modern era of data science is increasingly defined by datasets whose volume exceeds these capabilities, entering the realm of "Big Data." This paradigm requires a fundamental shift in our computational architecture, from single-node processing to distributed computing, and a corresponding shift in our algorithmic thinking.

\section{Associative Statistics: The Foundation of Scalability}
\label{sec:associative_statistics}
The key to processing data in a distributed, parallel manner is the ability to break a large computation into smaller, independent pieces, and then combine the partial results to obtain the final, global result. An algorithm is scalable in this manner only if the statistic it computes is \textbf{associative}.

\subsection*{Formal Definition of an Associative Statistic}
A statistic $s(D)$ is associative if, for any disjoint partition of the data $D = D_1 \cup D_2 \cup \dots \cup D_r$, the global statistic $s(D)$ can be perfectly reconstructed by applying a simple combining function $g$ to the partial statistics computed on each partition. For a two-set partition $D = D_1 \cup D_2$:
$$ s(D_1 \cup D_2) = g(s(D_1), s(D_2)) $$
The crucial property is that the result is independent of how the data was partitioned. For example, the sum is associative with the combining function being addition itself. This allows us to compute the sum of a massive dataset by summing chunks in parallel and then summing the partial sums. The count is similarly associative. Since the mean is simply the sum divided by the count, it can also be computed scalably. In contrast, statistics like the median are not associative; there is no general way to combine the medians of two subsets to find the median of the whole.

\subsection*{The Augmented Moment Matrix for Multivariate Statistics}
This concept extends to multivariate statistics. To scalably compute the mean vector and covariance matrix of a high-dimensional dataset, we need an associative statistic that captures all the necessary information. This is the \textbf{augmented moment matrix}.

For each data vector $\mathbf{x}_i \in \mathbb{R}^p$, we create an augmented vector $\mathbf{w}_i$ by prepending a 1: $\mathbf{w}_i = [1, \mathbf{x}_i^T]^T$. The augmented moment matrix $\mathbf{A}$ is the sum of the outer products of these augmented vectors:
$$ \mathbf{A} = \sum_{i=1}^{n} \mathbf{w}_i \mathbf{w}_i^T = 
\begin{pmatrix}
n & \sum \mathbf{x}_i^T \\
\sum \mathbf{x}_i & \sum \mathbf{x}_i \mathbf{x}_i^T
\end{pmatrix}
=
\begin{pmatrix}
n & n\mathbf{m}^T \\
n\mathbf{m} & \mathbf{M}
\end{pmatrix}
$$
where $\mathbf{m}$ is the mean vector and $\mathbf{M}$ is the matrix of second moments. Because this matrix is formed by a simple sum, it is an associative statistic. We can compute $\mathbf{A}_j$ for each partition $D_j$ on a separate machine, and then compute the global matrix as $\mathbf{A} = \sum_j \mathbf{A}_j$. From the final global matrix $\mathbf{A}$, we can directly extract the total count $n$, the mean vector $\mathbf{m}$, and the matrix of second moments $\mathbf{M}$, from which the covariance matrix can be easily calculated: $\mathbf{\Sigma} = \frac{1}{n}\mathbf{M} - \mathbf{m}\mathbf{m}^T$.

\section{The MapReduce Paradigm}
\label{sec:mapreduce}
\textbf{MapReduce} is a programming model that formalizes and abstracts the process of distributed computation using associative statistics. It allows data scientists to specify the logic of a large-scale computation without needing to manage the complexities of parallelization, fault tolerance, and data distribution. It consists of three conceptual stages:

\begin{enumerate}
    \item \textbf{Map Stage}: A user-defined "mapper" function is applied in parallel to every chunk of the raw input data. Its job is to process each record and emit a set of intermediate (key, value) pairs. The choice of key is critical, as it defines the axis of aggregation. For example, in a word count problem, the mapper would emit `(word, 1)` for every word it encounters.
    \item \textbf{Shuffle Stage}: This is an automatic stage managed by the framework. It collects all emitted key-value pairs from all mappers and groups them by key, such that all values associated with the same key are routed to a single worker node.
    \item \textbf{Reduce Stage}: A user-defined "reducer" function is applied to the list of values for each key. Its job is to aggregate, summarize, or transform this list into a final result. For the word count example, the reducer for the key `'hello'` would receive a list `[1, 1, 1, ...]` and would sum it to produce the final output `('hello', total_count)`.
\end{enumerate}
Frameworks like \textbf{Apache Hadoop}, with its Hadoop Distributed File System (HDFS) for redundant storage and YARN for resource management, provide the robust, fault-tolerant infrastructure required to execute MapReduce jobs at massive scale.